\section{Adjoint functors}\label{sec:adjoint_functors}

\paragraph{Comma categories}

\begin{definition}\label{def:diagonal_functor}\mimprovised
  For a fixed category \( \cat{C} \), we define the \( \cat{I} \)-indexed \term[en=diagonal functor (\cite[64]{MacLane1998CategoryTheory})]{diagonal functor} \( {\Delta_{\cat{C}}^{\cat{I}}: \cat{C} \to [\cat{I}, \cat{C}]} \) as the \hyperref[con:currying]{uncurrying} of the \hyperref[def:product_category/projection]{projection functor} \( \Pi_{\cat{C}}: \cat{C} \times \cat{I} \to \cat{C} \).

  This functor sends every object \( A \) of \( \cat{C} \) to the \term{constant functor}
  \begin{equation*}
    \begin{aligned}
      &\Delta_{\cat{C}}^{\cat{I}}(A): \cat{I} \to \cat{C}, \\
      &\Delta_{\cat{C}}^{\cat{I}}(A)(i) \coloneqq A, \\
      &\Delta_{\cat{C}}^{\cat{I}}(A)(p: i \to j) \coloneqq \id_A,
    \end{aligned}
  \end{equation*}
  and every morphism \( f: A \to B \) to the natural transformation
  \begin{equation*}
    \Delta_{\cat{C}}^{\cat{I}}(f) \coloneqq \seq{ f }_{k \in \obj(\cat{I})}.
  \end{equation*}
\end{definition}
\begin{comments}
  \item The functor \( \Delta_{\cat{C}}^{\cat{I}} \) (without the subscript and superscript) is called a \enquote{diagonal functor} in \cite[64]{MacLane1998CategoryTheory} and a \enquote{constant functor} in \cite[194]{Awodey2010CategoryTheory}. Its definition is a projection is due to Awodey.

  \item For \hyperref[def:comma_category]{comma categories}, we will usually take \( \cat{I} \) to be the \hyperref[def:singleton_category]{terminal category} \( \cat{1} \). As long as \( \cat{C} \) is a small category, \( \Delta_{\cat{C}}^{\cat{1}}(A) \) is simply a \hyperref[def:global_element]{global element} of \( \cat{C} \) in \( \cat{Cat} \) corresponding to \( A \).

  \item It is called a diagonal functor because, if \( \cat{I} \) is a discrete category with two objects, then the functor \( \Delta_{\cat{C}}^{\cat{I}}: \cat{C} \to [\cat{I}, \cat{C}] \) describes the same information as the functor \( D: \cat{C} \to \cat{C} \times \cat{C} \) that sends an object \( A \) to the ordered pair \( (A, A) \).
\end{comments}

\begin{definition}\label{def:comma_category}\mcite[45; 46]{MacLane1998CategoryTheory}
  Suppose we have the \hyperref[def:functor]{functors}
  \begin{equation*}
    \begin{aligned}
      \includegraphics[page=1]{output/def__comma_category}
    \end{aligned}
  \end{equation*}

  We define the \term{comma category} \( (F \downarrow G) \), with notes on special important cases where either \( F \) or \( G \) is a \hyperref[def:diagonal_functor]{constant functor} --- either \( F = \Delta_{\cat{C}}^{\cat{1}}(A) \) or \( G = \Delta_{\cat{C}}^{\cat{1}}(B) \).

  \begin{thmenum}
    \thmitem{def:comma_category/objects} The \hyperref[def:category/objects]{objects} of \( (F \downarrow G) \) are triples \( (A, B, s) \), where \( A \) and \( B \) are objects of \( \cat{A} \) and \( \cat{B} \), respectively, and \( s: F(A) \to G(B) \).

    We may omit \( A \) if \( F = \Delta_{\cat{C}}^{\cat{1}}(A) \) --- a pair \( (B, s) \), where \( s: A \to G(B) \), suffices. Similarly, we may omit \( B \) if \( G = \Delta_{\cat{C}}^{\cat{1}}(B) \).

    \thmitem{def:comma_category/morphisms} The \hyperref[def:category/morphisms]{morphisms} between \( (A, B, s) \) and \( (A', B', s') \) are pairs \( (f, g) \), where \( f: A \to A' \) and \( g: B \to B' \), such that the following diagram commutes:
    \begin{equation*}
      \begin{aligned}
        \includegraphics[page=2]{output/def__comma_category}
      \end{aligned}
    \end{equation*}

    When \( F = \Delta_{\cat{C}}^{\cat{1}}(A) \) or \( G = \Delta_{\cat{C}}^{\cat{1}}(B) \), respectively, this diagram simplifies to
    \begin{equation*}
      \begin{aligned}
        \includegraphics[page=3]{output/def__comma_category}
        \qquad
        \includegraphics[page=4]{output/def__comma_category}
      \end{aligned}
    \end{equation*}

    In the first case, \( f \) is necessarily \( \id_A \) and can be omitted, while in the second case \( g = \id_B \).

    \thmitem{def:comma_category/composition} The \hyperref[def:category/composition]{composition} of two morphisms is their componentwise composition, like in \hyperref[def:product_category]{product categories}.

    \thmitem{def:comma_category/identity} The \hyperref[def:category/identity]{identity morphism} on the object \( (A, s, B) \) is the identity pair \( (\id_A, \id_B) \).
  \end{thmenum}
\end{definition}
\begin{comments}
  \item The degenerate comma category \( (\Delta_{\cat{C}}^{\cat{I}}(A), \Delta_{\cat{C}}^{\cat{I}}(B)) \) is \hyperref[def:isomorphism_of_categories]{isomorphic} to the \hyperref[def:discrete_category]{discrete category} on \( \cat{C}(A, B) \) --- its objects are triples \( (A, B, s) \), where \( s: A \to B \), and its morphisms are simply pairs \( (\id_A, \id_B) \) of identity morphisms.

  \Textcite[47]{MacLane1998CategoryTheory} explains that the notation \enquote{\( \hom(A, B) \)} is the reason for the name \enquote{comma categories}, as well as for the notation \( (F, G) \) --- \enquote{a notation which we avoid because the comma is already overworked}.
\end{comments}

\paragraph{Universal morphisms}

\begin{definition}\label{def:concrete_category}\mcite[def. 5.1]{AdámekHerrlichStrecker2004JoyOfCats}
  A \term[en=concrete category (\cite[def. 5.1]{AdámekHerrlichStrecker2004JoyOfCats})]{relatively concrete category} over the \term{base category} \( \cat{X} \) is a pair \( (\cat{C}, U) \), where \( \cat{C} \) is a category and \( U: \cat{C} \to \cat{X} \) is a \hyperref[def:functor_invertibility/faithful]{faithful functor}.

  If \( \cat{X} \) is the category \hyperref[def:category_of_small_sets]{\( \cat{Set} \)} of \hyperref[def:small_set]{small sets}, we will simply call \( (\cat{C}, U) \) a \term[ru=конкретная категория (\cite[228]{ЦаленкоШульгейфер1974ОсновыТеорииКатегорий}), en=concrete category (\cite[26]{MacLane1998CategoryTheory})]{concrete category}. For every morphism \( f: A \to B \) in \( \cat{C} \), \( U(f) \) is a function between the (small) sets \( U(A) \) and \( U(B) \).
\end{definition}
\begin{comments}
  \item This definition of relatively concrete (with respect to \( \cat{X} \)) categories is taken from \cite[def. 5.1]{AdámekHerrlichStrecker2004JoyOfCats}, although the authors use \enquote{concrete} rather than \enquote{relatively concrete}, and call a category concrete over \( \cat{Set} \) a \enquote{construct}. The book is based around the more general relatively concrete categories.

  Usually, only concrete categories over \( \cat{Set} \) are considered, for example in \cite[26]{MacLane1998CategoryTheory} and \cite[def. 1.6.17]{Riehl2016CategoryTheory}. This is pushed even further in \cite[228]{ЦаленкоШульгейфер1974ОсновыТеорииКатегорий}, where \( U \) is assumed to be a full embedding into \( \cat{Set} \) and not merely a faithful functor.

  \item Cancellation and invertibility of morphisms in concrete categories are discussed in \cref{thm:concrete_category_function_invertibility}.
\end{comments}

\begin{concept}\label{con:forgetful_functor}
  Most often, in a \hyperref[def:concrete_category]{concrete category} \( (\cat{C}, U) \), the object \( U(A) \) \enquote{forgets} some structure present in \( A \). It is thus conventional to call \( U \) a \term[ru=пренебрегающий функтор (\cite[180]{ЦаленкоШульгейфер1974ОсновыТеорииКатегорий})]{forgetful functor} and any \hyperref[def:adjoint_functor]{left adjoint} to \( U \) a \term{free functor}.
\end{concept}

\begin{remark}\label{rem:concrete_category_notation}
  Outside this chapter, most \hyperref[def:concrete_category]{concrete categories} will come with an obvious \hyperref[con:forgetful_functor]{forgetful functor}, and we will avoid writing the functor explicitly.
\end{remark}

\begin{definition}\label{def:universal_morphism}\mcite[55]{MacLane1998CategoryTheory}
  \hyperref[def:comma_category]{Comma categories} allow formalizing the most common pattern of \hyperref[con:universal_property]{universal properties}.

  \begin{thmenum}
    \thmitem{def:universal_morphism/source} Fix a functor \( U: \cat{D} \to \cat{C} \), which we may regard as \hyperref[con:forgetful_functor]{forgetful} (even if \( U \) is not faithful). In the comma category \( (\Delta_{\cat{C}}^{\cat{1}} (A) \downarrow U) \), the universal property expressing that \( (\widetilde{A}, \iota) \) is an \hyperref[def:universal_objects/initial]{initial object} can be stated as follows:
    \begin{displayquote}
      For every object \( X \) in \( \cat{D} \) and every morphism \( {f: A \to U(X)} \) in \( \cat{C} \), there exists a unique morphism \( \widetilde{f}: \widetilde{A} \to X \) in \( \cat{D} \) such that the following diagram commutes:
      \begin{equation*}
        \begin{aligned}
          \includegraphics[page=1]{output/def__universal_morphism}
        \end{aligned}
      \end{equation*}
    \end{displayquote}

    Here \( \widetilde{A} \) can be regarded as \enquote{the closest approximation} of \( A \) in \( \cat{D} \), and \( \iota \) --- as an inclusion. By the above diagram, we can lift a morphism from \( A \) to an object of \( \cat{D} \) to a unique morphism in \( \cat{D} \). As shown in \cref{thm:adjunction_from_universal_morphism/unit}, this property can be used to characterize a \hyperref[def:adjoint_functor]{left adjoint functor}.

    \thmitem{def:universal_morphism/sink} Dually, in the category \( (U \downarrow \Delta_{\cat{C}}^{\cat{1}} (A)) \) (with the functors exchanged), the universal property expressing that \( (\widetilde{A}, \pi) \) is a \hyperref[def:universal_objects/terminal]{terminal object} can be stated as follows:
    \begin{displayquote}
      For every object \( X \) in \( \cat{D} \) and every morphism \( f: U(X) \to A \) in \( \cat{C} \), there exists a unique morphism \( \widetilde{f}: X \to \widetilde{A} \) in \( \cat{D} \) such that the following diagram commutes:
      \begin{equation*}
        \begin{aligned}
          \includegraphics[page=2]{output/def__universal_morphism}
        \end{aligned}
      \end{equation*}
    \end{displayquote}

    Here \( \widetilde{A} \) plays an analogous role, but is instead characterized by the morphisms \emph{into} \( A \). Instead of an inclusion \( \iota: A \to U(\widetilde{A}) \), here we have a projection \( \pi: U(\widetilde{A}) \to A \).

    As shown in \cref{thm:adjunction_from_universal_morphism/counit}, this property can be used to characterize a \hyperref[def:adjoint_functor]{right adjoint functor}.
  \end{thmenum}

  We refer to either \( (A, \iota) \) or \( (A, \pi) \) as \term[en=universal morphism (\cite[233]{MacLane1971CategoricalAlgebraFoundations}) / universal arrow (\cite[55]{MacLane1998CategoryTheory})]{universal morphisms}.
\end{definition}
\begin{comments}
  \item In accordance with \cref{rem:concrete_category_notation}, we usually find it convenient to skip writing the forgetful functor \( U \).
\end{comments}

\begin{remark}\label{rem:representable_functor_universal_morphism}
  \hyperref[con:universal_property]{Universal properties} can be described via \hyperref[def:functor_representation]{representable functors}.

  As in \cref{def:universal_morphism}, consider some functor \( U: \cat{C} \to \cat{D} \), which we may regard as \hyperref[con:forgetful_functor]{forgetful} (even if \( U \) is not faithful). Here we additionally need to assume that \( \cat{C} \) is \hyperref[def:small_category]{locally small}.

  For a fixed object \( A \) in \( \cat{C} \), consider the functor \( V \coloneqq H^A \bincirc U \):
  \begin{equation*}
    \begin{aligned}
      &V: \cat{C} \to \cat{Set}, \\
      &V(X) \coloneqq \cat{C}(A, U(X)), \\
      &V(g: X \to Y) \coloneqq (\underbrace{s}_{\mathclap{A \to U(X) \quad}} \mapsto \underbrace{U(g) \bincirc s}_{A \to U(Y)}). \\
    \end{aligned}
  \end{equation*}

  Suppose that \( (\widetilde{A}, \rho) \) is a \hyperref[def:functor_representation]{representation} of \( V \). By \cref{rem:yoneda_representation_characterization}, there exists a unique morphism \( \iota: A \to U(\widetilde{A}) \) such that \( \rho \) can be described by the natural transformation from \fullref{thm:yonedas_lemma}:
  \begin{equation*}
    \alpha_{V,\widetilde{A}}(\iota) \coloneqq \seq[\big]{ \underbrace{\widetilde{f}}_{\widetilde{A} \to X} \mapsto \underbrace{\overbrace{V(\widetilde{f})(\iota)}^{U(\widetilde{f}) \bincirc \iota}}_{A \mapsto U(X)} }_{X \in \cat{C}}.
  \end{equation*}

  Since \( \rho \) is by assumption an isomorphism, by \cref{rem:yoneda_representation_characterization}, for every object \( X \) and every morphism \( f: A \to U(X) \), there exists a unique morphism \( \widetilde{f}: \widetilde{A} \to X \) such that \( f = U(\widetilde{f}) \bincirc \iota \).

  This is precisely the universal property from \cref{def:universal_morphism/source} that characterizes \( (\widetilde{A}, \iota) \) as a universal morphism. In fact, \textcite[62]{Riehl2016CategoryTheory} defines a \enquote{universal property} as a representable functor \( V \) along with the \enquote{universal element} \( (\widetilde{A}, \iota) \) that uniquely determines a representation \( (\widetilde{A}, \rho) \).
\end{remark}

\begin{theorem}[Discrete category universal property]\label{thm:discrete_category_universal_property}
  Denote by \( U \) the \hyperref[con:forgetful_functor]{forgetful functor} which sends a \hyperref[def:category]{category} to its set of objects.

  The \hyperref[def:discrete_category]{discrete category} \( D(A) \) is the unique up to an \hyperref[def:isomorphism_of_categories]{isomorphism} category that satisfies the following \hyperref[con:universal_property]{universal property}:
  \begin{displayquote}
    For every category \( \cat{C} \) and every function \( f: A \to U(\cat{C}) \), there exists a unique \hyperref[def:functor]{functor} \( \widetilde{f}: F(A) \to \cat{C} \) such that the following diagram commutes:
    \begin{equation*}
      \begin{aligned}
        \includegraphics[page=1]{output/thm__discrete_category_universal_property}
      \end{aligned}
    \end{equation*}
  \end{displayquote}
\end{theorem}
\begin{comments}
  \item \Cref{thm:adjunction_from_universal_morphism/unit} implies that \( D \) is \hyperref[def:adjoint_functor]{left adjoint} to \( U \).
\end{comments}
\begin{proof}
  For every element \( a \) of \( A \), the functor \( \widetilde{f} \) must satisfy
  \begin{equation*}
    \widetilde{f}(a) = \widetilde{f}(\iota_A(a)) = f(a),
  \end{equation*}
  which enforces the definition
  \begin{equation*}
    \begin{aligned}
      &\widetilde{f}: F(A) \to \cat{C} \\
      &\widetilde{f}(a) \coloneqq f(a), \\
      &\widetilde{f}(\id_a) \coloneqq \id_{f(a)}.
    \end{aligned}
  \end{equation*}
\end{proof}

\begin{theorem}[Free category universal property]\label{thm:free_category_universal_property}
  Denote by \( U \) the \hyperref[con:forgetful_functor]{forgetful functor} which sends a \hyperref[def:category]{category} to its underlying \hyperref[def:directed_multigraph]{directed multigraph}.

  The \hyperref[def:free_category]{free category} \( F(G) \) is the unique up to an isomorphism category that satisfies the following \hyperref[con:universal_property]{universal property}:
  \begin{displayquote}
    For every category \( \cat{C} \) and every graph homomorphism \( \varphi: G \to U(\cat{C}) \), there exists a unique \hyperref[def:functor]{functor} \( \Phi: F(G) \to \cat{C} \) such that the following diagram commutes:
    \begin{equation*}
      \begin{aligned}
        \includegraphics[page=1]{output/thm__free_category_universal_property}
      \end{aligned}
    \end{equation*}
  \end{displayquote}
\end{theorem}
\begin{comments}
  \item \Cref{thm:adjunction_from_universal_morphism/unit} implies that \( F \) is \hyperref[def:adjoint_functor]{left adjoint} to \( U \).
\end{comments}
\begin{proof}
  For every node \( u \) in \( G \), the functor \( \Phi \) must satisfy
  \begin{equation*}
    \Phi(u) = \Phi((\iota_G)_V(u)) = \varphi_V(u).
  \end{equation*}

  Similarly, for every arc \( u \reloset{a} \to v \) in \( G \), the functor \( \Phi \) must satisfy
  \begin{equation*}
    \Phi(a) = \Phi((\iota_G)_A(a)) = \varphi_A(a).
  \end{equation*}

  These conditions give an unambiguous definition of \( \Phi \).
\end{proof}

\begin{definition}\label{def:factors_through}\mimprovised
  In some \hyperref[def:category]{category}, consider an object \( X \) and a morphism \( f: A \to B \). We say that \( f \) \term{factors through} \( X \) if there exist morphisms \( g: A \to X \) and \( h: X \to B \) such that the following diagram commutes:
  \begin{equation}\label{eq:def:factors_through}
    \begin{aligned}
      \includegraphics[page=1]{output/def__factors_through}
    \end{aligned}
  \end{equation}

  If \( g \) and \( h \) are uniquely determined by \( f \) and \( X \), we say that \( f \) \term{uniquely factors through} \( X \).
\end{definition}
\begin{comments}
  \item This concept is somehow \hyperref[con:assumed_knowledge]{assumed knowledge} and formal definitions are difficult to find.
\end{comments}

\begin{theorem}[Quotient category universal property]\label{thm:quotient_category_universal_property}
  For every category \( \cat{C} \) and every \hyperref[def:congruence_in_category]{congruence} \( {\cong} \) on \( \cat{C} \), the \hyperref[def:quotient_category]{quotient category} \( \cat{C} / {\cong} \) has the following \hyperref[con:universal_property]{universal property}:
  \begin{displayquote}
    Every functor \( F: \cat{C} \to \cat{D} \) for which \( f \cong f' \) implies \( F(f) = F(f') \) \hyperref[def:factors_through]{uniquely factors through} \( \cat{C} / {\cong} \).

    More precisely, there exists a unique functor \( \widetilde{F}: \cat{C} / {\cong} \to \cat{D} \), such that the following diagram commutes:
    \begin{equation*}
      \begin{aligned}
        \includegraphics[page=1]{output/thm__quotient_category_universal_property}
      \end{aligned}
    \end{equation*}
  \end{displayquote}
\end{theorem}
\begin{proof}
  For any morphism \( f \) in \( \cat{C} \), we want
  \begin{equation*}
    \widetilde{F}(\pi(f)) = F(f).
  \end{equation*}

  Because of our restriction on \( F \), \( \pi(f) = \pi(f') \) implies \( F(f) = F(f') \), so the above can be used as a definition.

  Uniqueness holds by construction.
\end{proof}

\paragraph{Adjunctions}

\begin{remark}\label{rem:adjunction_of_functors}
  Adjoint functors help systemize \hyperref[def:universal_morphism]{universal morphisms}.

  \begin{thmenum}
    \thmitem{rem:adjunction_of_functors/hom} The concept is attributed by \textcite[79]{MacLane1998CategoryTheory} to \textcite[def. 3.1]{Kan1958AdjointFunctors}, who, given oppositely-directed functors \( F: \cat{C} \to \cat{D} \) and \( G: \cat{D} \to \cat{C} \) and a \hyperref[thm:natural_isomorphism]{natural isomorphism} with components
    \begin{equation*}
      \alpha_{A,X}: \cat{D}(F(A), X) \to \cat{C}(A, G(X)),
    \end{equation*}
    says that that \enquote{\( F \) is left adjoint to \( G \) under \( \alpha \)} and that \enquote{\( G \) is right adjoint to \( F \) under \( \alpha \)}.

    Thus, \( \alpha \) is a natural transformation between the following twisted variants of the \hyperref[def:hom_functor]{\( \hom \)-bifunctor}\fnote{The second functor generalizes the functor discussed in \cref{rem:representable_functor_universal_morphism}}:
    \begin{paracol}{2}
      \begin{leftcolumn}
        \ParacolAlignmentHack
        \begin{equation*}
          \begin{aligned}
            &L: \cat{C}^\oppos \times \cat{D} \to \cat{Set}, \\
            &L(A, X) \coloneqq \cat{D}(F(A), X), \\
            &L(\underbrace{f}_{B \to A}, \underbrace{g}_{X \to Y}) \coloneqq \parens[\big]{ \underbrace{s}_{\mathclap{F(A) \to X}} \mapsto \underbrace{g \bincirc s \bincirc F(f)}_{\mathclap{F(B) \to Y}} },
          \end{aligned}
        \end{equation*}
      \end{leftcolumn}

      \begin{rightcolumn}
        \ParacolAlignmentHack
        \begin{equation*}
          \begin{aligned}
            &R: \cat{C}^\oppos \times \cat{D} \to \cat{Set}, \\
            &R(A, X) \coloneqq \cat{C}(A, G(X)), \\
            &R(\underbrace{f}_{B \to A}, \underbrace{g}_{X \to Y}) \coloneqq \parens[\big]{ \underbrace{s}_{\mathclap{A \to G(X)}} \mapsto \underbrace{G(g) \bincirc s \bincirc f}_{\mathclap{B \to G(Y)}} },
          \end{aligned}
        \end{equation*}
      \end{rightcolumn}
    \end{paracol}

    This definition is later refined to call the triple \( (F, G, \alpha) \) an \enquote{adjunction}, like we will do in \cref{def:hom_adjunction}. Such a definition is given in
    \cite[80]{MacLane1998CategoryTheory},
    \cite[def. 16.4.1]{Schubert1972Categories},
    \cite[def. 4.1.1]{Perrone2024StartingCategoryTheory},
    \cite[def. 2.1.1]{Leinster2014BasicCategoryTheory},
    \cite[116]{Riehl2016CategoryTheory} and
    \cite[178]{ЦаленкоШульгейфер1974ОсновыТеорииКатегорий}.

    Formally, both \( \cat{C} \) and \( \cat{D} \) must be \hyperref[def:small_category]{locally small}, although the formalisms can be suitably generalized. Such a generalization is implicitly assumed by all the aforementioned authors. This particular issue is discussed in \cite{MathSE:definition_of_adjoint_functor_and_locally_small_categories}.

    \thmitem{rem:adjunction_of_functors/unit_counit} An alternative definition that works more generally uses two natural isomorphisms, the \enquote{unit} \( \eta: \id_{\cat{C}} \to G \bincirc F \) and \enquote{counit} \( \varepsilon: F \bincirc G \to \id_{\cat{D}} \), that must satisfy the conditions in \cref{def:unit_counit_adjunction}. \Textcite[def. 9.1]{Awodey2010CategoryTheory} gives a variation of this definition --- a natural transformation \( \eta: \id_{\cat{C}} \to G \bincirc F \) such that \( \eta_A \) is a universal morphism for every \( A \).

    Awodey also uses the symbol \enquote{\( U \)} instead of \enquote{\( G \)}, highlighting how \( G \) is often a \hyperref[con:forgetful_functor]{forgetful functor}.

    \thmitem{rem:adjunction_of_functors/order} The role of \( F \) and \( G \) in an adjunction is asymmetric. Oftentimes \( G \) is forgetful, while \( F \) generates \hyperref[con:free_construction]{free constructions} in \( \cat{D} \).

    Different authors may prefer different letters for denoting the categories \( \cat{C} \) and \( \cat{D} \) (e.g. \( \mscrA \) and \( \mscrB \)), but the alphabetic order that we use is usually followed. All cited authors use this convention except for Schubert, who prefers the opposite, where \( \cat{C} \) and \( \cat{D} \) are swapped and the free constructions belong to \( \cat{C} \).
  \end{thmenum}
\end{remark}

\begin{definition}\label{def:hom_adjunction}\mcite[def. 4.1.1]{Perrone2024StartingCategoryTheory}
  Suppose that \( \cat{C} \) and \( \cat{D} \) are \hyperref[def:small_category]{locally small} categories. A \term[en=adjunction (\cite[def. 4.1.1]{Perrone2024StartingCategoryTheory})]{\( \hom \)-adjunction} from the \hyperref[def:functor]{functor} \( F: \cat{C} \to \cat{D} \) to \( G: \cat{D} \to \cat{C} \) is a \hyperref[thm:natural_isomorphism]{natural isomorphism} \( \alpha \) with components
  \begin{equation}\label{eq:def:hom_adjunction}
    \alpha_{A,X}: \cat{D}(F(A), X) \Rightarrow \cat{C}(A, G(X)).
  \end{equation}
\end{definition}
\begin{comments}
  \item The two functors for which \( \alpha \) is a transformation are written out in \cref{rem:adjunction_of_functors/hom}.
\end{comments}

\begin{definition}\label{def:unit_counit_adjunction}\mcite[thm. 4.2.17]{Perrone2024StartingCategoryTheory}
  A \term[en=adjunction (\cite[def. 4.1.1]{Perrone2024StartingCategoryTheory})]{unit-counit adjunction} from the \hyperref[def:functor]{functor} \( F: \cat{C} \to \cat{D} \) to \( G: \cat{D} \to \cat{C} \) is a pair of \hyperref[thm:natural_isomorphism]{natural transformations}
  \begin{align*}
           \eta &: \id_{\cat{C}} \Rightarrow G \bincirc F, \\
    \varepsilon &: F \bincirc G \Rightarrow \id_{\cat{D}},
  \end{align*}
  called the \term[en=unit (\cite[def. 4.1.1]{Perrone2024StartingCategoryTheory})]{unit} and \term[en=counit (\cite[def. 4.1.11]{Perrone2024StartingCategoryTheory})]{counit}, respectively.

  We require the following triangle diagrams to commute for every object \( A \) in \( \cat{A} \) and \( X \) in \( \cat{D} \):
  \begin{equation}\label{eq:def:unit_counit_adjunction/triangle}
    \begin{aligned}
      \includegraphics[page=1]{output/def__unit_counit_adjunction}
      \quad
      \includegraphics[page=2]{output/def__unit_counit_adjunction}
    \end{aligned}
  \end{equation}
\end{definition}
\begin{comments}
  \item The signatures of \( \eta \) and \( \varepsilon \) are adjusted to coincide with the natural isomorphisms in the definition of equivalence in \cref{def:equivalence_of_categories}.
\end{comments}

\begin{definition}\label{def:adjoint_functor}\mimprovised
  We say that the functor \( F: \cat{C} \to \cat{D} \) is \term[en=left adjoint (\cite[79]{MacLane1998CategoryTheory})]{left adjoint} to \( G: \cat{D} \to \cat{C} \) and that \( G \) is \term[en=right adjoint (\cite[79]{MacLane1998CategoryTheory})]{right adjoint} to \( F \) if there exists a \hyperref[def:unit_counit_adjunction]{unit-counit adjunction} between them.

  If both \( \cat{C} \) and \( \cat{D} \) are \hyperref[def:small_category]{locally small}, \cref{thm:unit_counit_adjunction_induces_hom_adjunction} constructs an explicit \hyperref[def:hom_adjunction]{\( \hom \)-adjunction} from this unit-counit adjunction. Conversely, \cref{thm:hom_adjunction_induces_unit_counit_adjunction} constructs a unit-counit adjunction from a \( \hom \)-adjunction.

  We will usually use \cref{thm:adjunction_from_universal_morphism} to characterize an adjunction via \hyperref[def:universal_morphism]{universal morphisms}.
\end{definition}
\begin{comments}
  \item We take the unit-counit definition as primitive due to size considerations. The definition is discussed in more detail in \cref{rem:adjunction_of_functors}.
\end{comments}

\begin{proposition}\label{thm:hom_adjunction_induces_unit_counit_adjunction}
  For every \hyperref[def:hom_adjunction]{\( \hom \)-adjunction} \( \alpha \) from \( F: \cat{C} \to \cat{D} \) to \( G: \cat{D} \to \cat{C} \), the indexed families
  \begin{align*}
    \eta        &\coloneqq \seq[\big]{ \alpha_{A,F(A)}(\id_{F(A)}) }_{A \in \obj(\cat{C})}, \\
    \varepsilon &\coloneqq \seq[\big]{ \alpha_{G(X),X}^{-1}(\id_{G(X)}) }_{X \in \obj(\cat{D})}
  \end{align*}
  are natural transformations that form a \hyperref[def:unit_counit_adjunction]{unit-counit adjunction}.
\end{proposition}
\begin{proof}
  \SubProof{Proof of naturality of \( \eta \)} Fix a morphism \( t: A \to B \) in \( \cat{C} \). By the naturality of \( \alpha \), the following diagram commutes:
  \begin{equation*}
    \begin{aligned}
      \includegraphics[page=1]{output/thm__hom_adjunction_induces_unit_counit_adjunction}
    \end{aligned}
  \end{equation*}

  Thus,
  \begin{equation*}
    \cat{C}(\id_A, [G \bincirc F](t))(\underbrace{ \varphi_{A, F(A)}(\id_{F(A)}) }_{\eta_A}) = \varphi_{A, F(B)}(\underbrace{ \cat{D}(F(\id_A), F(t))(\id_{F(A)}) }_{F(t) \bincirc \id_{F(A)}})
  \end{equation*}
  and
  \begin{equation*}
    \cat{C}(t, G(\id_{F(B)}))(\underbrace{ \varphi_{B, F(B)}(\id_{F(B)}) }_{\eta_B}) = \varphi_{A, F(B)}(\underbrace{ \cat{D}(F(t), \id_{F(B)})(\id_{F(B)}) }_{\id_{F(B)} \bincirc F(t)}).
  \end{equation*}

  The right sides of the above equalities coincide, so
  \begin{equation*}
    [G \bincirc F](t) \bincirc \eta_A
    =
    \cat{C}(\id_A, [G \bincirc F](t))(\eta_A)
    =
    \cat{C}(t, G(\id_{F(B)}))(\eta_B)
    =
    \eta_B \bincirc t.
  \end{equation*}

  This equality is encoded in the following diagram:
  \begin{equation*}
    \begin{aligned}
      \includegraphics[page=2]{output/thm__hom_adjunction_induces_unit_counit_adjunction}
    \end{aligned}
  \end{equation*}

  This shows the naturality of \( \eta \).

  \SubProof{Proof of naturality of \( \varepsilon \)} Analogously, for any morphism \( p: X \to Y \) in \( \cat{C} \), the following diagram commutes:
  \begin{equation*}
    \begin{aligned}
      \includegraphics[page=3]{output/thm__hom_adjunction_induces_unit_counit_adjunction}
    \end{aligned}
  \end{equation*}

  Thus,
  \begin{equation*}
    \cat{D}(F(\id_{G(X)}), p)(\underbrace{ \varphi^{-1}_{G(X), X}(\id_{G(X)}) }_{\varepsilon_X}) = \varphi^{-1}_{G(X), Y}(\underbrace{ \cat{C}(\id_{G(X)}, G(p))(\id_{G(X)}) }_{G(p) \bincirc \id_{G(X)}})
  \end{equation*}
  and
  \begin{equation*}
    \cat{D}([F \bincirc G](p), \id_Y)(\underbrace{ \varphi^{-1}_{G(Y), Y}(\id_{G(Y)}) }_{\varepsilon_Y}) = \varphi^{-1}_{G(X), Y}(\underbrace{ \cat{C}(G(p), G(\id_Y))(\id_{G(Y)}) }_{\id_{G(Y)} \bincirc G(p)}).
  \end{equation*}

  The right sides are again equal, so the following diagram commutes:
  \begin{equation*}
    \begin{aligned}
      \includegraphics[page=4]{output/thm__hom_adjunction_induces_unit_counit_adjunction}
    \end{aligned}
  \end{equation*}

  \SubProof{Proof of commutativity of triangle for \( G(X) \)} By the naturality of \( \alpha \), the following diagram commutes:
  \begin{equation*}
    \begin{aligned}
      \includegraphics[page=5]{output/thm__hom_adjunction_induces_unit_counit_adjunction}
    \end{aligned}
  \end{equation*}

  Starting in the upper-left corner with the morphism \( \id_{[F \bincirc G](X)} \), we obtain
  \begin{equation*}
    \underbrace{\cat{C}(\id_{G(X)}, G(\varepsilon_X))\parens[\big]{ \overbrace{\alpha_{G(X), [F \bincirc G](X)}(\id_{[F \bincirc G](X)})}^{\eta_{G(X)}} }}_{G(\varepsilon_X) \bincirc \eta_{G(X)}}
    =
    \underbrace{\alpha_{G(X), X} \parens[\big]{ \overbrace{\cat{D}(F(\id_{G(X)}), \varepsilon_X)(\id_{[F \bincirc G](X)})}^{\id_{\varepsilon_X \bincirc [F \bincirc G](X)}} }}_{\id_{G(X)}}.
  \end{equation*}

  \SubProof{Proof of commutativity of triangle for \( F(A) \)} Similarly, the following diagram commutes:
  \begin{equation*}
    \begin{aligned}
      \includegraphics[page=6]{output/thm__hom_adjunction_induces_unit_counit_adjunction}
    \end{aligned}
  \end{equation*}

  Starting in the lower-right corner with the morphism \( \id_{[G \bincirc F](A)} \), we obtain
  \begin{equation*}
    \underbrace{\cat{D}(F(\eta_A), \id_{F(A)})\parens[\big]{ \overbrace{\alpha_{[G \bincirc F](A), F(A)}^{-1} (\id_{[G \bincirc F](A)})}^{\varepsilon^{F(A)}} }}_{\varepsilon_{F(A)} \bincirc F(\eta_A)}
    =
    \underbrace{\alpha_{A, F(A)}^{-1} \parens[\big]{ \overbrace{\cat{C}(\eta_A, G(\id_{F(A)}))(\id_{[G \bincirc F](A)})}^{\id_{[G \bincirc F](A)} \bincirc \eta_A} }}_{\id_{F(A)}}.
  \end{equation*}
\end{proof}

\begin{proposition}\label{thm:unit_counit_adjunction_induces_hom_adjunction}
  Assuming \( \cat{C} \) and \( \cat{D} \) are \hyperref[def:small_category]{locally small}, for every \hyperref[def:unit_counit_adjunction]{unit-counit adjunction} \( (\eta, \varepsilon) \), the indexed family
  \begin{equation*}
    \alpha \coloneqq \seq[\big]{ \underbrace{s}_{\mathclap{F(A) \to X}} \mapsto \underbrace{G(s) \bincirc \eta_A}_{A \to G(X)} }_{A \in \obj(\cat{C}), X \in \obj(\cat{D})}, \\
  \end{equation*}
  is a \hyperref[def:hom_adjunction]{\( \hom \)-adjunction}, i.e. a natural isomorphism between the twisted \( \hom \)-functors from \cref{rem:adjunction_of_functors/hom}. Its inverse is
  \begin{equation*}
    \beta \coloneqq \seq[\big]{ \underbrace{t}_{\mathclap{A \to G(X)}} \mapsto \underbrace{\varepsilon_X \bincirc F(t)}_{F(A) \to X} }_{A \in \obj(\cat{C}), X \in \obj(\cat{D})}
  \end{equation*}
\end{proposition}
\begin{proof}
  \SubProof{Proof of naturality of \( \alpha \)} Given morphisms \( f: B \to A \) in \( \cat{C} \) and \( g: X \to Y \) in \( \cat{D} \), we must show that the following diagram commutes:
  \begin{equation*}
    \begin{aligned}
      \includegraphics[page=1]{output/thm__unit_counit_adjunction_induces_hom_adjunction}
    \end{aligned}
  \end{equation*}

  Fix a morphism \( s: F(A) \to X \). On one side, we have
  \begin{equation*}
    \alpha_{B,Y}\parens[\big]{ \cat{D}(F(f), g)(s) } = \alpha_{B,Y}(g \bincirc s \bincirc F(f)) = G(g \bincirc s \bincirc F(f)) \bincirc \eta_B.
  \end{equation*}

  On the other side, we have
  \begin{equation*}
    \cat{C}(f, G(g))\parens[\big]{ \alpha_{A,X}(s) } = \cat{C}(f, G(g))\parens[\big]{ G(s) \bincirc \eta_A } = G(g) \bincirc G(s) \bincirc \eta_A \bincirc f.
  \end{equation*}

  It remains to show that \( \eta_A \bincirc f = [G \bincirc F](f) \bincirc \eta_B \), which readily follows from the naturality of \( \eta \):
  \begin{equation*}
    \begin{aligned}
      \includegraphics[page=2]{output/thm__unit_counit_adjunction_induces_hom_adjunction}
    \end{aligned}
  \end{equation*}

  \SubProof{Proof that \( \beta \) is the inverse of \( \alpha \)} Fix a morphism \( s: F(A) \to G \) in \( \cat{D} \). First, we must show that \( s \) equals
  \begin{equation*}
    [\beta_{A,X} \bincirc \alpha_{A,X}](s) = \varepsilon_X \bincirc F(G(s) \bincirc \eta_A) = \varepsilon_X \bincirc [F \bincirc G](s) \bincirc F(\eta_A).
  \end{equation*}

  This equality follows by combining the triangle diagram for \( F(A) \) with the naturality diagram for \( s \):
  \begin{equation*}
    \begin{aligned}
      \includegraphics[page=3]{output/thm__unit_counit_adjunction_induces_hom_adjunction}
    \end{aligned}
  \end{equation*}

  Thus, \( \beta_{A,X} \) is a left inverse of \( \alpha_{A,X} \).

  Conversely, every morphism \( t: A \to G(X) \) in \( \cat{C} \) must equal
  \begin{equation*}
    [\alpha_{A,X} \bincirc \beta_{A,X}](t) = G(\varepsilon_X \bincirc F(t)) \bincirc \eta_A = G(\varepsilon_X) \bincirc [G \bincirc F](t) \bincirc \eta_A,
  \end{equation*}
  which follows from the diagram
  \begin{equation*}
    \begin{aligned}
      \includegraphics[page=4]{output/thm__unit_counit_adjunction_induces_hom_adjunction}
    \end{aligned}
  \end{equation*}

  Thus, \( \beta_{A,X} \) is also a right inverse of \( \alpha_{A,X} \).

  We have already shown that \( \alpha \) is a natural transformation, so \cref{thm:natural_isomorphism} implies that \( \beta \) is also a natural transformation and, in fact, an inverse of \( \alpha \).
\end{proof}

\begin{proposition}\label{thm:opposite_unit_counit_adjunction}
  The functor \( F: \cat{C} \to \cat{D} \) is \hyperref[def:adjoint_functor]{left adjoint} to \( G: \cat{D} \to \cat{C} \) if and only if the \hyperref[def:opposite_functor]{poosite functor} \( F^\oppos \) is right adjoint to \( G^\oppos \).

  Consider oppositely-directed functors \( F: \cat{C} \to \cat{D} \) and \( G: \cat{D} \to \cat{C} \). The pair of natural transformations
  \begin{align*}
           \eta &: \id_{\cat{C}} \Rightarrow G \bincirc F, \\
    \varepsilon &: F \bincirc G \Rightarrow \id_{\cat{D}},
  \end{align*}
  forms a \hyperref[def:unit_counit_adjunction]{unit-counit adjunction} if and only if the \hyperref[def:opposite_natural_transformation]{opposite natural transformations}
  \begin{align*}
    \varepsilon^\oppos &: \id_{\cat{D}} \Rightarrow F^\oppos \bincirc G^\oppos, \\
           \eta^\oppos &: G^\oppos \bincirc F^\oppos \Rightarrow \id_{\cat{C}}
  \end{align*}
  form an adjunction from the \hyperref[def:opposite_functor]{opposite functor} \( G^\oppos \) to \( F^\oppos \).
\end{proposition}
\begin{comments}
  \item \Cref{thm:composition_of_opposite_functors} implies that \( (F \bincirc G)^\oppos = F^\oppos \bincirc G^\oppos \) and \( (G \bincirc F)^\oppos = G^\oppos \bincirc F^\oppos \).

  \item This is part of the categorical duality principles listed in \cref{ex:thm:category_duality}.
\end{comments}
\begin{proof}
  Straightforward.
\end{proof}

\begin{corollary}\label{thm:opposites_of_adjoint_functors}
  The functor \( F \) is \hyperref[def:adjoint_functor]{left adjoint} to \( G \) if and only if the \hyperref[def:opposite_functor]{opposite functor} \( F^\oppos \) is right adjoint to \( G^\oppos \).
\end{corollary}
\begin{proof}
  Follows from \cref{thm:opposite_unit_counit_adjunction}.
\end{proof}

\begin{proposition}\label{thm:adjunction_from_universal_morphism}
  We will show how \hyperref[def:universal_morphism]{universal morphisms} are related to \hyperref[def:unit_counit_adjunction]{unit-counit adjunctions}. Consider oppositely-directed functors \( F: \cat{C} \to \cat{D} \) and \( G: \cat{D} \to \cat{C} \).

  \begin{thmenum}
    \thmitem{thm:adjunction_from_universal_morphism/unit} The natural transformation
    \begin{equation*}
      \eta: \id_{\cat{C}} \Rightarrow G \bincirc F.
    \end{equation*}
    is the unit of an adjunction if and only if, for every object \( A \) in \( \cat{C} \), the pair \( (A, \eta_A) \) is a universal morphism, i.e. if it satisfies the following universal property:
    \begin{displayquote}
      For every object \( X \) in \( \cat{D} \) and every morphism \( {f: A \to G(X)} \) in \( \cat{C} \), there exists a unique morphism \( \widetilde{f}: F(A) \to X \) in \( \cat{D} \) such that the following diagram commutes:
      \begin{equation*}
        \begin{aligned}
          \includegraphics[page=1]{output/thm__adjunction_from_universal_morphism}
        \end{aligned}
      \end{equation*}
    \end{displayquote}

    When \( G \) is treated as forgetful, this property can be used to characterize \( F \) as a left adjoint of \( G \).

    \thmitem{thm:adjunction_from_universal_morphism/counit} Dually, the natural transformation
    \begin{equation*}
      \varphi: F \bincirc G \Rightarrow \id_{\cat{D}}
    \end{equation*}
    is a counit of an adjunction if and only if, for every object \( X \) in \( \cat{D} \), the pair \( (X, \varepsilon_X) \) is a universal morphism, i.e. if it satisfies the following universal property:
    \begin{displayquote}
      For every object \( A \) in \( \cat{C} \) and every morphism \( {f: F(A) \to X} \) in \( \cat{D} \), there exists a unique morphism \( \widetilde{f}: A \to G(X) \) in \( \cat{C} \) such that the following diagram commutes:
      \begin{equation*}
        \begin{aligned}
          \includegraphics[page=2]{output/thm__adjunction_from_universal_morphism}
        \end{aligned}
      \end{equation*}
    \end{displayquote}

    When \( F \) is treated as forgetful, this property can be used to characterize \( G \) as a right adjoint of \( F \).
  \end{thmenum}
\end{proposition}
\begin{comments}
  \item In accordance with \cref{rem:concrete_category_notation}, we usually find it convenient to skip writing the forgetful functor \( U \). Furthermore, the naturality of \( \eta \) and \( \varepsilon \) is often trivial and requires no verification. We discuss these issues in \cref{ex:def:adjoint_functor}.
\end{comments}
\begin{proof}
  \SubProofOf{thm:adjunction_from_universal_morphism/unit}
  \SufficiencySubProof* Suppose that \( (\eta, \varepsilon) \) is a unit-counit adjunction.

  Fix some object \( X \) in \( \cat{D} \) and some morphism \( f: A \to G(X) \). To determine \( \widetilde{f} \), suppose that it satisfies its defining universal property. This can be combined with the universality of \( \eta \) and the triangle diagram for \( F(A) \) into the diagram
  \begin{equation*}
    \begin{aligned}
      \includegraphics[page=3]{output/thm__adjunction_from_universal_morphism}
    \end{aligned}
  \end{equation*}

  This leads to the definition
  \begin{equation*}
    \widetilde{f} \coloneqq F(f) \bincirc \eta_{G(X)},
  \end{equation*}
  and incidentally proves the uniqueness of \( \widetilde{f} \).

  It remains to show that \( G(\widetilde{f}) \bincirc \eta_A = f \), which follows by combining the triangle diagram for \( G(X) \) with the naturality diagram for \( f \):
  \begin{equation*}
    \begin{aligned}
      \includegraphics[page=4]{output/thm__adjunction_from_universal_morphism}
    \end{aligned}
  \end{equation*}

  Therefore, \( (A, \eta_A) \) is a universal morphism.

  \NecessitySubProof* Suppose that, for every object \( A \) in \( \cat{C} \), the pair \( (A, \eta_A) \) is a universal morphism.

  In particular, for every object \( X \) in \( \cat{D} \), the following diagram commutes:
  \begin{equation*}
    \begin{aligned}
      \includegraphics[page=5]{output/thm__adjunction_from_universal_morphism}
    \end{aligned}
  \end{equation*}

  We define the counit transformation \( \varepsilon \) as by assigning \( \widetilde{\id_{G(X)}} \) to the component \( \varepsilon_X \). Thus, the triangle diagram for \( G(X) \) is satisfied by definition.

  This triangle diagram can then be utilized to prove naturality of \( \varepsilon \):
  \begin{equation*}
    \begin{aligned}
      \includegraphics[page=6]{output/thm__adjunction_from_universal_morphism}
    \end{aligned}
  \end{equation*}

  To conclude that \( (\eta, \varepsilon) \) is an adjunction, we need to show that the triangle diagram for \( F(X) \) commutes. This follows from the universal property because, by assumption, \( {\varepsilon_{F(A)} = \widetilde{G(\id_{F(A)})}} \) is the unique morphism in \( \cat{D} \) such that the following diagram commutes:
  \begin{equation*}
    \begin{aligned}
      \includegraphics[page=7]{output/thm__adjunction_from_universal_morphism}
    \end{aligned}
  \end{equation*}

  Because of this uniqueness, we can conclude that the triangle diagram for \( F(A) \) also commutes:
  \begin{equation*}
    \begin{aligned}
      \includegraphics[page=8]{output/thm__adjunction_from_universal_morphism}
    \end{aligned}
  \end{equation*}

  \SubProofOf{thm:adjunction_from_universal_morphism/counit} \Cref{thm:opposite_unit_counit_adjunction} implies that \( (\eta, \varepsilon) \) is an adjunction from \( F \) to \( G \) if and only if \( (\varepsilon^\oppos, \eta^\oppos) \) is an adjunction from \( G^\oppos \) to \( F^\oppos \).

  For the latter, \cref{thm:adjunction_from_universal_morphism/unit} implies that \( \varepsilon^\oppos \) is the unit of an adjunction if and only if, for every object \( X \) in \( \cat{D}^\oppos \), the pair \( (X, \varepsilon_X^*) \) satisfies the following universal property:
  \begin{displayquote}
    For every object \( A \) in \( \cat{C}^\oppos \) and every morphism \( {f^*: X \to F^\oppos(A)} \) in \( \cat{D}^\oppos \), there exists a unique morphism \( (\widetilde{f^*})^*: G^\oppos(X) \to A \) in \( \cat{C}^\oppos \) such that the following diagram commutes:
    \begin{equation*}
      \begin{aligned}
        \includegraphics[page=9]{output/thm__adjunction_from_universal_morphism}
      \end{aligned}
    \end{equation*}
  \end{displayquote}

  Reversing the morphisms where appropriate, we obtain the following universal property:
  \begin{displayquote}
    For every object \( A \) in \( \cat{C} \) and every morphism \( {f: F(X) \to A} \) in \( \cat{D} \), there exists a unique morphism \( \widetilde{f^*}: A \to G(X) \) in \( \cat{C} \) such that the following diagram commutes:
    \begin{equation*}
      \begin{aligned}
        \includegraphics[page=10]{output/thm__adjunction_from_universal_morphism}
      \end{aligned}
    \end{equation*}
  \end{displayquote}

  Therefore, \( (\eta, \varepsilon) \) is an adjunction if and only if the above holds for every pair \( (X, \eta_X) \).
\end{proof}

\begin{corollary}\label{thm:functor_adjoint_uniqueness}
  If a functor has two \hyperref[def:adjoint_functor]{left adjoints} (resp. right adjoints), then there exists a unique \hyperref[def:natural_isomorphism]{natural isomorphism} between them.
\end{corollary}
\begin{proof}
  Fix a functor \( G: \cat{D} \to \cat{C} \) and suppose that both \( F: \cat{D} \to \cat{C} \) and \( F': \cat{D} \to \cat{D} \) are left adjoint, with adjunctions \( (\eta, \varepsilon) \) and \( (\eta', \varepsilon') \).

  \Cref{thm:adjunction_from_universal_morphism} implies that, for every object \( A \) in \( \cat{C} \), both \( (F(A), \varepsilon_A) \) and \( (F'(A), \varepsilon_A') \) are \hyperref[def:universal_morphism]{universal morphisms}, i.e. \hyperref[def:universal_objects/initial]{initial objects} in the \hyperref[def:comma_category]{comma category} \( (\Delta_{\cat{C}}^{\cat{1}} (A) \downarrow G) \).

  \Cref{thm:def:universal_objects/unique} implies that \( (F(A), \varepsilon_A) \) and \( (F'(A), \varepsilon_A') \) are isomorphic. Then there exists a unique isomorphism \( \varphi_A: F(A) \to F'(A) \) such that the following diagram commutes:
  \begin{equation*}
    \begin{aligned}
      \includegraphics[page=1]{output/thm__functor_adjoint_uniqueness}
    \end{aligned}
  \end{equation*}

  These isomorphisms form a natural isomorphism \( \varphi: F \to F' \). Indeed, all its components are by definition invertible and unique, while naturality can be shown as follows. Fix a morphism \( f: A \to B \). The naturality of \( \eta \) allows us to construct the commutative diagram:
  \begin{equation*}
    \begin{aligned}
      \includegraphics[page=2]{output/thm__functor_adjoint_uniqueness}
    \end{aligned}
  \end{equation*}

  Due to the uniqueness of \( \varphi_A \) and \( \varphi_A' \), it follows that the following diagram, expressing naturality of \( \varphi \), commutes:
  \begin{equation*}
    \begin{aligned}
      \includegraphics[page=3]{output/thm__functor_adjoint_uniqueness}
    \end{aligned}
  \end{equation*}

  Therefore, \( \varphi: F \to F' \) is the unique isomorphism in \( [F, F'] \).
\end{proof}

\begin{example}\label{ex:def:adjoint_functor}
  We list examples of \hyperref[def:adjoint_functor]{adjoint functors}:
  \begin{thmenum}
    \thmitem{ex:def:adjoint_functor/free_construction} The free constructions, discussed in \cref{con:free_construction}, follow the pattern from \cref{thm:adjunction_from_universal_morphism/unit}.

    For categories, two examples of free constructions are \fullref{thm:discrete_category_universal_property} and \fullref{thm:free_category_universal_property}, which differ by the objects we construct a category from.

    For example, in the latter, we start with a \hyperref[def:directed_multigraph]{directed multigraph} \( G \) and its \hyperref[def:free_category]{free category} \( F(G) \), with the obvious inclusion \( \iota_G: G \to [U \bincirc F](G) \). The universal property of \( (G, \iota_G) \) is as follows:
    \begin{displayquote}
      For every category \( \cat{C} \) and every graph homomorphism \( \varphi: G \to U(\cat{C}) \), there exists a unique \hyperref[def:functor]{functor} \( \Phi: F(G) \to \cat{C} \) such that the following diagram commutes:
      \begin{equation*}
        \begin{aligned}
          \includegraphics[page=1]{output/ex__def__adjoint_functor}
        \end{aligned}
      \end{equation*}
    \end{displayquote}

    \thmitem{ex:def:adjoint_functor/us_ds}\mcite{MathOF:free_digraph} Denote the \hyperref[def:directed_graph/category]{category of  small directed graphs} by \( \cat{DGph} \) and the \hyperref[def:directed_graph/category]{category of small undirected graphs} by \( \cat{UGph} \).

    The doubling functor \( D_S: \cat{UGph} \to \cat{DGph} \) from \cref{def:graph_functors/simple_doubling} is a natural candidate for left adjoint to the forgetful functor \hyperref[def:graph_functors/simple_doubling]{\( U_S: \cat{DGph} \to \cat{UGph} \)}. This is not the case, however, as there are \hyperref[def:undirected_graph/morphism]{undirected graph homomorphisms} that fail to be \hyperref[def:directed_graph/morphism]{directed graph homomorphisms}. An example is shown in \cref{fig:ex:def:adjoint_functor/us_ds}.

    \begin{figure}
      \hfill
      \includegraphics[page=2]{output/ex__def__adjoint_functor}
      \hfill
      \includegraphics[page=3]{output/ex__def__adjoint_functor}
      \hfill\hfill
      \caption{Two undirected homomorphisms from \( G \) to \( U_S(Q) \), denoted using dashed lines, only one of which is a directed multigraph homomorphism from \( D_S(G) \) to \( Q \).}
      \label{fig:ex:def:adjoint_functor/us_ds}
    \end{figure}

    If we consider instead a directed graph homomorphisms \( f: Q \to D_S(G) \), it is clear that it is also an undirected graph homomorphism from \( U_S(Q) \) to \( G \). Thus, \( D_S \) is actually \emph{right adjoint} to \( U_S \).
  \end{thmenum}
\end{example}

\begin{proposition}\label{thm:concrete_category_function_invertibility}
  Fix a \hyperref[def:concrete_category]{concrete category} \( \cat{C} \) with forgetful functor \( U: \cat{C} \to \cat{Set} \).

  \begin{thmenum}
    \thmitem{thm:concrete_category_function_invertibility/injective} Every \hyperref[def:function_invertibility/injective]{injective} morphism in \( \cat{C} \) is a \hyperref[def:morphism_cancellation/left]{monomorphism}.

    The converse holds if \( U \) has a \hyperref[def:adjoint_functor]{left adjoint}, or if the monomorphism splits.

    \thmitem{thm:concrete_category_function_invertibility/surjective} Dually, every \hyperref[def:function_invertibility/surjective]{surjective} morphism in \( \cat{C} \) is an \hyperref[def:morphism_cancellation/right]{epimorphism}.

    The converse holds if \( U \) has a \hyperref[def:adjoint_functor]{right adjoint}, or if the epimorphism splits.

    \thmitem{thm:concrete_category_function_invertibility/bijective} A \hyperref[def:function_invertibility/bijective]{bijective} morphism \( \cat{C} \) is a \hyperref[def:morphism_invertibility/isomorphism]{categorical isomorphism} if and only if its \hyperref[def:set_valued_map/inverse]{set-theoretic inverse} is also a morphism in \( \cat{C} \).
  \end{thmenum}
\end{proposition}
\begin{proof}
  \SubProofOf{thm:concrete_category_function_invertibility/injective} Consider the morphism \( f: A \to B \).

  \SufficiencySubProof* Suppose that the function \( U(f) \) is injective.

  \Cref{thm:function_invertibility_categorical/left_cancellable} implies that \( U(f) \) is left cancellable on functions. \Cref{thm:def:functor_invertibility/faithful_reflects_cancellation} implies that \( f \) is also left cancellable.

  \SubProof*{Proof of necessity if \( U \) has a left adjoint} Suppose that \( g: X \to Y \) is a monomorphism in \( \cat{C} \).

  Let \( F: \cat{Set} \to \cat{C} \) be a left adjoint to \( U \). Fix functions \( {f_1: A \to U(X)} \) and \( {f_2: A \to U(X)} \) such that
  \begin{equation*}
    U(g) \bincirc f_1 = U(g) \bincirc f_2.
  \end{equation*}

  By \cref{thm:adjunction_from_universal_morphism/unit}, there exists a unique morphism \( \widetilde{f_1}: F(A) \to X \) in \( \cat{C} \) such that the following diagram commutes:
  \begin{equation*}
    \begin{aligned}
      \includegraphics[page=1]{output/thm__concrete_category_function_invertibility}
    \end{aligned}
  \end{equation*}

  We have
  \begin{equation*}
    U(g) \bincirc f_1 = U(g) \bincirc U(\widetilde{f_1}) \bincirc \eta_A = U(g \bincirc \widetilde{f_1}) \bincirc \eta_A,
  \end{equation*}
  thus \( h \coloneqq g \bincirc \widetilde{f_1} \) is the unique morphism in \( \cat{C} \) such that the following diagram commutes:
  \begin{equation*}
    \begin{aligned}
      \includegraphics[page=2]{output/thm__concrete_category_function_invertibility}
    \end{aligned}
  \end{equation*}

  But \( U(g) \bincirc f_1 = U(g) \bincirc f_2 \) by assumption, so we can replace \( f_1 \) with \( f_2 \) in the above diagram and conclude that \( h = g \bincirc \widetilde{f_2} \).

  Since \( g \) is a monomorphism, it follows that \( \widetilde{f_1} = \widetilde{f_2} \). By the uniqueness of \( \widetilde{f_1} \), it follows that \( f_1 = f_2 \).

  We conclude that \( U(f) \) is also cancellable. \Cref{thm:function_invertibility_categorical/left_cancellable} implies that \( U(f) \) is injective.

  \SubProof*{Proof of necessity if \( f \) is left invertible} Suppose that \( f: A \to B \) is left invertible, i.e. a split monomorphism.

  If \( U(f) \) is empty, it is injective as a consequence of \cref{thm:def:functor/two_sided_inverses}.

  Otherwise, \( U(f) \) is injective as a consequence of \cref{thm:function_invertibility_categorical/nonempty_left_invertible}.

  \SubProofOf{thm:concrete_category_function_invertibility/surjective} Follows by duality.

  \SubProofOf{thm:concrete_category_function_invertibility/bijective} Consider the morphism \( f: A \to B \).

  \SufficiencySubProof* Suppose that \( U(f) \) is bijective and that there exists a morphism \( g: B \to A \) such that \( U(g) \) is the two-sided inverse of \( U(f) \).

  Then
  \begin{equation*}
    \underbrace{\id_{U(A)}}_{U(\id_A)} = \underbrace{U(g) \bincirc U(f)}_{U(g \bincirc f)}.
  \end{equation*}

  Since \( U \) is faithful, we conclude that \( g \bincirc f = \id_A \).

  We can analogously prove that \( f \bincirc g = \id_B \). Therefore, \( g \) is the two-sided inverse of \( f \), i.e. \( f \) is an isomorphism.

  \NecessitySubProof* Conversely, suppose that \( f: A \to B \) is an isomorphism.

  \Cref{thm:def:functor/two_sided_inverses} implies that \( U(f^{-1}) \) is the two-sided inverse of \( U(f) \), i.e. \( U(f) \) is bijective.
\end{proof}

\paragraph{Adjoint equivalences}

\begin{definition}\label{def:adjoint_equivalence}\mcite[93]{MacLane1998CategoryTheory}
  We call the \hyperref[def:equivalence_of_categories]{equivalence} \( (F, G, \eta, \varepsilon) \) an \term{adjoint equivalence} if \( (\eta, \varepsilon) \) is a \hyperref[def:unit_counit_adjunction]{unit-counit adjunction}.
\end{definition}

\begin{proposition}\label{thm:adjoint_equivalence}
  Let \( (F, G, \eta, \varepsilon) \) be a \hyperref[def:equivalence_of_categories]{equivalence} between \( \cat{C} \) and \( \cat{D} \).

  Then there exists a \hyperref[def:natural_isomorphism]{natural isomorphism} \( \zeta: \id_{\cat{C}} \Rightarrow G \bincirc F \) such that \( (F, G, \zeta, \varepsilon) \) is an \hyperref[def:adjoint_equivalence]{adjoint equivalence}.
\end{proposition}
\begin{proof}
  From \fullref{thm:categorical_equivalence_characterization} it follows that \( F \) is fully faithful and essentially surjective.

  Then, for every object \( A \) in \( \cat{A} \), it has a partial inverse
  \begin{equation*}
    M_A: \cat{D}\parens[\big]{ F(A), [F \bincirc G \bincirc F](A) } \to \cat{C}\parens[\big]{ A, [F \bincirc G](A) }.
  \end{equation*}

  Hence, we can define
  \begin{equation*}
    \zeta_A \coloneqq M(\varepsilon_{F(A)}^{-1})
  \end{equation*}
  so that
  \begin{equation*}
    F(\zeta_A) = F(M_A(\varepsilon_{F(A)}^{-1})) = \varepsilon_{F(A)}^{-1}.
  \end{equation*}

  The family of morphisms \( \zeta \) is a natural transformation because \( \varepsilon^{-1} \) is. Indeed, given a morphism \( f: A \to B \) in \( \cat{C} \), the following diagram commutes:
  \begin{equation*}
    \begin{aligned}
      \includegraphics[page=1]{output/thm__adjoint_equivalence}
    \end{aligned}
  \end{equation*}

  Taking the preimages under \( F \), as justified by \cref{thm:def:functor_invertibility/faithful_reflects_composition}, we obtain the naturality diagram for \( f \):
  \begin{equation*}
    \begin{aligned}
      \includegraphics[page=2]{output/thm__adjoint_equivalence}
    \end{aligned}
  \end{equation*}

  Furthermore, by \cref{thm:def:functor_invertibility/fully_faithful_reflects_invertible}, every component \( \zeta_A \) is an isomorphism, so \cref{thm:natural_isomorphism} implies that \( \zeta \) is a natural isomorphism.

  It remains to show that \( (\zeta, \varepsilon) \) is an adjunction. The triangle diagram for \( F(A) \) is simple enough:
  \begin{equation*}
    \begin{aligned}
      \includegraphics[page=3]{output/thm__adjoint_equivalence}
    \end{aligned}
  \end{equation*}

  Similarly, the following diagram also commutes:
  \begin{equation*}
    \begin{aligned}
      \includegraphics[page=4]{output/thm__adjoint_equivalence}
    \end{aligned}
  \end{equation*}

  Taking the preimages under \( F \), we obtain the triangle diagram for \( G(X) \):
  \begin{equation*}
    \begin{aligned}
      \includegraphics[page=5]{output/thm__adjoint_equivalence}
    \end{aligned}
  \end{equation*}
\end{proof}
